%%%%%%%%%%%%%%%%%%%%%%%%%%%%%%%%%%%%%%%%%%%%%%%%%%%%%%%%%%%%%%%%%%%%%%%%%%%%%%%%
\begin{frame}[fragile]\frametitle{}
\begin{center}
{\Large Agent Patterns in LangGraph}
\end{center}
\end{frame}
%%%%%%%%%%%%%%%%%%%%%%%%%%%%%%%%%%%%%%%%%%%%%%%%%%%%%%%%%%%%%%%%%%%%%%%%%%%%%%%%
\begin{frame}[fragile]\frametitle{Introduction to Agents and Runtime}
      \begin{itemize}
        \item Agent: LLM-powered action taker that performs reasoning and tool invocation.
        \item Agent runtime: Loop that manages agent execution and interactions.
        \item In LangChain: Runtime handled by \texttt{AgentExecutor} running a fixed execution loop.
        \item Tools and reasoning executed sequentially in a single loop.
        \item In LangGraph: Runtime redesigned as a flexible graph-based executor.
        \item Supports cycles, branching, and dynamic decision-making.
        \item Each node represents a logical step (tool call, reasoning, or condition).
        \item Enables more control and reusability across complex workflows.
      \end{itemize}
\end{frame}
%%%%%%%%%%%%%%%%%%%%%%%%%%%%%%%%%%%%%%%%%%%%%%%%%%%%%%%%%%%%%%%%%%%%%%%%%%%%%%%%
\begin{frame}[fragile]\frametitle{Agentic Patterns: Autonomous vs Workflow Automation}
      \begin{itemize}
        \item \textbf{Autonomous Agents}:
        \begin{itemize}
            \item High-level goals with minimal constraints
            \item Self-directed exploration and decision-making
            \item Example: "Research this topic and write a report"
            \item Higher risk of unpredictability
        \end{itemize}
        \item \textbf{Workflow Automation Agents}:
        \begin{itemize}
            \item Structured processes with defined checkpoints
            \item Predictable paths with conditional logic
            \item Example: "Review email → Classify → Route → Draft response"
            \item Balance between automation and control
        \end{itemize}
        \item \textbf{LangGraph Sweet Spot}: Workflow automation with controlled autonomy
        \item Graphs enable explicit workflow definition while allowing intelligent decisions at each node
      \end{itemize}
\end{frame}
%%%%%%%%%%%%%%%%%%%%%%%%%%%%%%%%%%%%%%%%%%%%%%%%%%%%%%%%%%%%%%%%%%%%%%%%%%%%%%%%
\begin{frame}[fragile]\frametitle{}
\begin{center}
{\Large Reflection Agents in LangGraph}

{\tiny (Ref: LangGraph Crash Course - Harish Neel)}

\end{center}
\end{frame}
%%%%%%%%%%%%%%%%%%%%%%%%%%%%%%%%%%%%%%%%%%%%%%%%%%%%%%%%%%%%%%%%%%%%%%%%%%%%%%%%
\begin{frame}[fragile]\frametitle{Understanding Reflection - The Human Perspective}
      \begin{itemize}
        \item Reflection means looking at yourself or your actions, like viewing yourself in a mirror
        \item Self-reflection after giving a presentation - thinking about how it went
        \item Re-reading an email after writing to ensure clarity and correctness
        \item Considering if a decision was the right choice after making it
        \item Having a mirror in front of you to think about improvements
        \item Thinking about what you've done in the past and how to do better
        \item Foundation concept for understanding AI reflection patterns
      \end{itemize}
\end{frame}
%%%%%%%%%%%%%%%%%%%%%%%%%%%%%%%%%%%%%%%%%%%%%%%%%%%%%%%%%%%%%%%%%%%%%%%%%%%%%%%%
\begin{frame}[fragile]\frametitle{What is a Reflection Agent Pattern?}
      \begin{itemize}
        \item AI system pattern that can look at its own outputs and improve them
        \item Similar to humans looking at themselves in a mirror and self-reflecting
        \item Makes AI systems better through iterative self-improvement
        \item Basic reflection agent system consists of two core components
        \item Generator agent - creates initial output or response
        \item Reflector agent - analyzes and critiques the generated output
        \item Both agents work together in a collaborative feedback loop
        \item Enables continuous improvement through multiple iterations
      \end{itemize}
\end{frame}
%%%%%%%%%%%%%%%%%%%%%%%%%%%%%%%%%%%%%%%%%%%%%%%%%%%%%%%%%%%%%%%%%%%%%%%%%%%%%%%%
\begin{frame}[fragile]\frametitle{Practical Example: Viral Tweet Generator}
      \begin{itemize}
        \item Simple application demonstrating basic reflection agent pattern
        \item Start node initiates the process with user input topic
        \item Tweet Generation Agent creates initial tweet based on provided topic
        \item Conditional edge routes output to Tweet Critiquing Agent
        \item Tweet Critiquing Agent acts as critical reviewer of generated content
        \item Critique focuses on viral potential: length, CTAs, hooks, hashtags
        \item Feedback loop continues for 4-6 iterations for improvement
        \item Each iteration produces progressively better, more viral-worthy tweets
        \item Real-world application potential as product feature for customers
      \end{itemize}
\end{frame}
%%%%%%%%%%%%%%%%%%%%%%%%%%%%%%%%%%%%%%%%%%%%%%%%%%%%%%%%%%%%%%%%%%%%%%%%%%%%%%%%
\begin{frame}[fragile]\frametitle{System 1 vs System 2: Traditional vs Reflection Patterns}
      \begin{itemize}
        \item \textbf{System 1 (Traditional):} One-shot prompt approach like ChatGPT/Claude
        \item Fast execution with immediate results
        \item Subconscious processing without active deep thinking
        \item Automatic response generation for everyday decisions
        \item More error-prone due to lack of reflection
        \item \textbf{System 2 (Reflection Pattern):} Two-agent collaborative approach
        \item Slower execution due to iterative feedback loops
        \item More conscious and effortful processing
        \item Better for complex decisions requiring precision
        \item Higher reliability through active thought-out responses
      \end{itemize}
\end{frame}
%%%%%%%%%%%%%%%%%%%%%%%%%%%%%%%%%%%%%%%%%%%%%%%%%%%%%%%%%%%%%%%%%%%%%%%%%%%%%%%%
\begin{frame}[fragile]\frametitle{Reflection Loop Process Flow}
      \begin{itemize}
        \item Initial generation phase creates first response/output
        \item Reflection phase analyzes generated content critically
        \item Reflector identifies merits: what went right in the output
        \item Reflector identifies issues: what went wrong and needs improvement
        \item Critique feedback sent back to generator brain/agent
        \item Generator incorporates feedback to revise and improve output
        \item Loop repeats for predetermined number of iterations (n times)
        \item Final improved output delivered to user after iteration completion
        \item Process applicable beyond tweets to any content generation task
      \end{itemize}
\end{frame}
%%%%%%%%%%%%%%%%%%%%%%%%%%%%%%%%%%%%%%%%%%%%%%%%%%%%%%%%%%%%%%%%%%%%%%%%%%%%%%%%
\begin{frame}[fragile]\frametitle{Three Types of Reflection Agents in LangGraph}
      \begin{itemize}
        \item \textbf{Basic Reflection Agents:} Foundation pattern with generator and reflector
        \item Reflector chain where reflector agent and generator agent work together
        \item Simple feedback loop for iterative improvement
        \item \textbf{Reflexion Agents:} Advanced pattern building on basic reflection
        \item Enhanced capabilities beyond basic reflection mechanisms
        \item \textbf{Language Agent Tree Search (LATS):} Most sophisticated approach
        \item Tree-based search methodology for complex decision making
        \item Recommended to master basic reflection first before advancing
        \item Progressive learning path from basic to advanced patterns
      \end{itemize}
\end{frame}
%%%%%%%%%%%%%%%%%%%%%%%%%%%%%%%%%%%%%%%%%%%%%%%%%%%%%%%%%%%%%%%%%%%%%%%%%%%%%%%%
\begin{frame}[fragile]\frametitle{Basic Reflection Agents}

Adding to graph

{\tiny (Ref: LangGraph Crash Course - Harish Neel)}

% [code-font: forced \scriptsize | overflow at the default size]
\begin{lstlisting}[basicstyle=\scriptsize\ttfamily]
from typing import List, Sequence
from langchain_core.messages import BaseMessage, HumanMessage
from langgraph.graph import END, START, StateGraph
from typing_extensions import TypedDict, Annotated
from langgraph.graph.message import add_messages
from chains import generation_chain, reflection_chain

# Define the state for our graph
class AgentState(TypedDict):
    messages: Annotated[Sequence[BaseMessage], add_messages]

def generate_node(state: AgentState):
    return {"messages": [generation_chain.invoke({"messages": state["messages"]})]}

def reflect_node(state: AgentState):
    response = reflection_chain.invoke({"messages": state["messages"]})
    return {"messages": [HumanMessage(content=response.content)]}

graph_builder = StateGraph(AgentState)
graph_builder.add_node("generate", generate_node)
graph_builder.add_node("reflect", reflect_node)
graph_builder.add_edge(START, "generate")
\end{lstlisting}
\end{frame}
%%%%%%%%%%%%%%%%%%%%%%%%%%%%%%%%%%%%%%%%%%%%%%%%%%%%%%%%%%%%%%%%%%%%%%%%%%%%%%%%
\begin{frame}[fragile]\frametitle{Basic Reflection Agents}

Adding to graph

{\tiny (Ref: LangGraph Crash Course - Harish Neel)}

% [code-font: forced \scriptsize | overflow at the default size]
\begin{lstlisting}[basicstyle=\scriptsize\ttfamily]
def should_continue(state: AgentState):
    if len(state['messages']) >= 6:
        # End after 3 generations (human + 3 x generate + 2 x reflect = 6)
        return END
    return "reflect"

graph_builder.add_conditional_edges("generate", should_continue)
graph_builder.add_edge("reflect", "generate")

app = graph_builder.compile()

# Visualize the graph
print(app.get_graph().draw_ascii())

# Invoke the graph
response = app.invoke({"messages": [HumanMessage(content="AI Agents taking over content creation")]})

print(response['messages'][-1].content)
\end{lstlisting}
\end{frame}
%%%%%%%%%%%%%%%%%%%%%%%%%%%%%%%%%%%%%%%%%%%%%%%%%%%%%%%%%%%%%%%%%%%%%%%%%%%%%%%%
\begin{frame}[fragile]\frametitle{}
\begin{center}
{\Large Structured LLM Outputs}
\end{center}
\end{frame}
%%%%%%%%%%%%%%%%%%%%%%%%%%%%%%%%%%%%%%%%%%%%%%%%%%%%%%%%%%%%%%%%%%%%%%%%%%%%%%%%
\begin{frame}[fragile]\frametitle{Introduction to Structured Outputs}
      \begin{itemize}
	\item LLMs traditionally return unstructured string responses that are difficult to process programmatically
	\item Structured outputs allow models to return data matching specific schemas we define
	\item Essential for software engineering where we work with objects, JSON, and database-ready formats
	\item Instead of random strings, get organized data with exact properties for manipulation
	\item Multiple output formats supported: JSON, dictionary, string, YAML, HTML
	\item Example: Request a joke and receive setup, punchline, and rating in structured JSON format
	\item Enables seamless integration with databases and application workflows
	\item Critical foundation for building advanced agentic AI systems like reflexion agents
	\item Transforms LLM responses from text blobs into actionable data structures
	\item Makes LLM outputs suitable for automated processing and storage systems
	  \end{itemize}
\end{frame}
%%%%%%%%%%%%%%%%%%%%%%%%%%%%%%%%%%%%%%%%%%%%%%%%%%%%%%%%%%%%%%%%%%%%%%%%%%%%%%%%
\begin{frame}[fragile]\frametitle{Pydantic Models for Structured Output}
      \begin{itemize}
	\item Pydantic is a Python library that defines data structures as blueprints for data validation
	\item Uses Python type hints (string, integer) to enforce correct data types automatically
	\item Define classes with required fields and descriptions for each property
	\item Integrated with LangChain using \texttt{with\_structured\_output()} method for LLM binding
	\item Works with any chat model (OpenAI, Groq, Llama) by simply swapping the model
	\item Internally converts Pydantic model into a tool that LLM must use exclusively
	\item Automatic validation ensures all properties are present with correct data types
	\item Throws errors automatically if validation fails, enabling fallback strategies
	\item LLM receives tool schema with property descriptions and required field specifications
	\item Response parsing and conversion handled automatically by LangChain framework
	  \end{itemize}
\end{frame}
%%%%%%%%%%%%%%%%%%%%%%%%%%%%%%%%%%%%%%%%%%%%%%%%%%%%%%%%%%%%%%%%%%%%%%%%%%%%%%%%
\begin{frame}[fragile]\frametitle{Pydantic Model Implementation Example}
      \begin{itemize}
	\item Use \texttt{with\_structured\_output()} to bind model to LLM instance
	\item Tool choice parameter forces LLM to use only the specified schema
	\item Request payload includes tools array with function type and parameter specifications
	\item Response contains structured JSON matching exact schema requirements
	\item LangChain handles parsing, validation, and Pydantic model instantiation automatically
	\item Validation failures trigger automatic error handling for robust applications
	  \end{itemize}

% [code-font: forced \scriptsize | overflow at the default size]
\begin{lstlisting}[language=Python, basicstyle=\scriptsize\ttfamily]
from pydantic import BaseModel, Field

class Country(BaseModel):
    """Information about a country"""
    name: str = Field(description="Name of the country")
    language: str = Field(description="Language of the country")
    capital: str = Field(description="Capital of the country")

structured_llm = model.with_structured_output(Country)
response = structured_llm.invoke("Tell me about France")
\end{lstlisting}
\end{frame}
%%%%%%%%%%%%%%%%%%%%%%%%%%%%%%%%%%%%%%%%%%%%%%%%%%%%%%%%%%%%%%%%%%%%%%%%%%%%%%%%
\begin{frame}[fragile]\frametitle{Modern Tool Binding with \texttt{.bind\_tools()}}
      \begin{itemize}
        \item While \texttt{with\_structured\_output} is good for forcing a single schema, the modern approach for general tool use is \texttt{.bind\_tools()}.
        \item This method attaches one or more tools (Pydantic models or functions) to an LLM.
        \item The model can then decide which tool to call, if any, based on the input prompt. This is the foundation of most agentic behavior.
      \end{itemize}
% [code-font: forced \tiny | overflow at the default size]
\begin{lstlisting}[language=Python, basicstyle=\tiny\ttfamily]
from langchain_groq import ChatGroq
from pydantic import BaseModel, Field

# Define a tool schema
class GetWeather(BaseModel):
    """Get the current weather in a given location"""
    location: str = Field(description="The city and state, e.g. San Francisco, CA")

# Initialize the model (any chat model with tool calling works)
llm = ChatGroq(model="openai/gpt-oss-20b")

# Bind the tool to the model
llm_with_tools = llm.bind_tools([GetWeather])

# Invoke the model
ai_msg = llm_with_tools.invoke("what is the weather in Boston?")

# The output contains a tool_calls attribute
print(ai_msg.tool_calls)
# > [{'name': 'GetWeather', 'args': {'location': 'Boston'}, 'id': '...'}]
\end{lstlisting}
\end{frame}
%%%%%%%%%%%%%%%%%%%%%%%%%%%%%%%%%%%%%%%%%%%%%%%%%%%%%%%%%%%%%%%%%%%%%%%%%%%%%%%%
\begin{frame}[fragile]\frametitle{Alternative Methods: TypedDict and JSON Schema}
      \begin{itemize}
	\item TypedDict approach uses Python typing without Pydantic validation overhead
	\item Annotated class provides type hints, default values, and property descriptions
	\item JSON schema method offers direct schema definition without Python classes
	\item Schema requires title, description, properties, and required field specifications
	  \end{itemize}

% [code-font: forced \tiny | overflow at the default size]
\begin{lstlisting}[language=Python, basicstyle=\tiny\ttfamily]
# TypedDict approach
from typing import TypedDict, Annotated, Optional

class JokeDict(TypedDict):
    setup: Annotated[str, "Setup of the joke"]
    punchline: Annotated[str, "Punchline to the joke"]
    rating: Annotated[Optional[int], None, "Joke rating 1-10"]

# JSON Schema approach
joke_schema = {"title": "Joke","type": "object",
    "properties": {
        "setup": {"type": "string", "description": "Setup of joke"},
        "punchline": {"type": "string", "description": "Punchline"}
    },"required": ["setup", "punchline"]
}
\end{lstlisting}
\end{frame}
%%%%%%%%%%%%%%%%%%%%%%%%%%%%%%%%%%%%%%%%%%%%%%%%%%%%%%%%%%%%%%%%%%%%%%%%%%%%%%%%
\begin{frame}[fragile]\frametitle{}
\begin{center}
{\Large Reflexion Agent in LangGraph}
\end{center}
\end{frame}
%%%%%%%%%%%%%%%%%%%%%%%%%%%%%%%%%%%%%%%%%%%%%%%%%%%%%%%%%%%%%%%%%%%%%%%%%%%%%%%%
\begin{frame}[fragile]\frametitle{Reflexion Agent System: Overview}
      \begin{itemize}
          \item Addresses drawbacks of basic reflection systems that lack grounding in live data
          \item Unlike reflection agents, reflexion agents fact-check with external data via API calls
          \item Combines self-critique with internet search capabilities for enriched content
          \item Generates citations and references from external sources during content creation
          \item Particularly useful for marketers writing blog posts requiring recent information
          \item Not limited to LLM training data cutoff dates like traditional reflection systems
          \item Iteratively improves content quality through multiple search-revise cycles
          \item Enables creation of well-researched, factually grounded content with live data
          \item Prevents hallucinations by validating information against current sources
          \item Supports episodic memory for context-aware, personalized interactions over time
      \end{itemize}
\end{frame}
%%%%%%%%%%%%%%%%%%%%%%%%%%%%%%%%%%%%%%%%%%%%%%%%%%%%%%%%%%%%%%%%%%%%%%%%%%%%%%%%
\begin{frame}[fragile]\frametitle{Reflexion Agent}

\begin{center}
\includegraphics[width=0.8\linewidth,keepaspectratio]{langgraph11}

{\tiny (Ref: LangGraph Crash Course - Harish Neel)}

\end{center}	  


\end{frame}
%%%%%%%%%%%%%%%%%%%%%%%%%%%%%%%%%%%%%%%%%%%%%%%%%%%%%%%%%%%%%%%%%%%%%%%%%%%%%%%%
\begin{frame}[fragile]\frametitle{Core Components Architecture}
      \begin{itemize}
          \item \textbf{Actor Agent:} Main driver that orchestrates all subcomponents and processes
          \item \textbf{Responder Agent:} Generates initial response with self-reflection capabilities
          \item \textbf{Revisor Agent:} Creates revised responses based on new external data
          \item \textbf{Tools/Tool Execution:} Handles API calls, primarily internet search functionality
          \item \textbf{Episodic Memory:} Recalls specific past interactions and experiences for context
          \item Actor contains all subcomponents: responder, revisor, and execute tools modules
          \item System operates through iterative loops of generation, search, and revision
          \item Each component has specific responsibilities within the overall workflow
          \item Tool execution primarily uses Tavily Search for gathering live internet data
          \item Architecture designed for modularity and extensibility with additional tools
      \end{itemize}
\end{frame}
%%%%%%%%%%%%%%%%%%%%%%%%%%%%%%%%%%%%%%%%%%%%%%%%%%%%%%%%%%%%%%%%%%%%%%%%%%%%%%%%
\begin{frame}[fragile]\frametitle{Workflow and Data Flow}
      \begin{itemize}
          \item User request goes to responder agent for initial content generation
          \item Responder outputs JSON with response, critique, and suggested search keywords
          \item Execute tools component searches internet using suggested keywords via Tavily
          \item Search results combined with initial response fed to revisor agent
          \item Revisor generates revised response with updated critique and citations
          \item Citations added to reference external sources used in content creation
          \item Process iterates: revised search terms trigger new searches and revisions
          \item Each iteration enriches content with more recent, relevant information
          \item Content length maintained (e.g., 250 words) while improving quality and accuracy
          \item Final output returned to user after multiple refinement cycles complete
      \end{itemize}
\end{frame}
%%%%%%%%%%%%%%%%%%%%%%%%%%%%%%%%%%%%%%%%%%%%%%%%%%%%%%%%%%%%%%%%%%%%%%%%%%%%%%%%
\begin{frame}[fragile]\frametitle{JSON Response Structure Example}
      \begin{itemize}
          \item Responder agent outputs structured JSON with three main properties
          \item \textbf{Response:} Contains the initial generated content (e.g., 250-word blog post)
          \item \textbf{Critique:} Self-assessment identifying areas for improvement in content
          \item \textbf{Search Terms:} Keywords suggested for gathering additional relevant information
          \item Revisor agent extends structure with additional citation property
          \item \textbf{Citations:} References to external sources with URLs and factual backing
          \item Revised versions update all properties: response, critique, search terms, citations
          \item Structure enables systematic improvement through multiple iterations
          \item JSON format facilitates programmatic processing and integration with LangGraph
          \item Supports grounding of content claims with verifiable external sources
      \end{itemize}
\end{frame}
%%%%%%%%%%%%%%%%%%%%%%%%%%%%%%%%%%%%%%%%%%%%%%%%%%%%%%%%%%%%%%%%%%%%%%%%%%%%%%%%
\begin{frame}[fragile]\frametitle{ReACT Architecture Introduction}
      \begin{itemize}
        \item ReACT (Reasoning and Acting) combines reasoning and acting capabilities
        \item Agent solves problems through continuous cycle: Reason \(\rightarrow\) Act \(\rightarrow\) Observe
        \item Flexible response to complex tasks using external tools
        \item Enhanced capabilities through tool integration
        \item LangChain 1.x provides a pre-built ReACT agent on top of LangGraph
        \item Tools can be search, image generation, databases, and more
        \item One call to \texttt{create\_agent} (it replaces the older \texttt{create\_react\_agent})
        \item Suitable for dynamic problem-solving scenarios
      \end{itemize}
\end{frame}
%%%%%%%%%%%%%%%%%%%%%%%%%%%%%%%%%%%%%%%%%%%%%%%%%%%%%%%%%%%%%%%%%%%%%%%%%%%%%%%%
\begin{frame}[fragile]\frametitle{ReACT Agent Implementation}
% [code-font: forced \scriptsize | overflow at the default size]
      \begin{lstlisting}[language=Python, basicstyle=\scriptsize\ttfamily]
import dotenv
from langchain.agents import create_agent
from langchain_core.tools import tool
from langchain_groq import ChatGroq

dotenv.load_dotenv()

# Tools are plain functions; the docstring is the description the model sees
@tool
def search_web(query: str) -> str:
    """Search the web for current events."""
    return f"Top results for: {query}"          # call a real search API here

@tool
def generate_image(prompt: str) -> str:
    """Generate an image from a text description and return its URL."""
    return "https://images.example.com/shark.png"   # call an image API here

tools = [search_web, generate_image]
model = ChatGroq(model="openai/gpt-oss-20b", temperature=0)

# Create ReACT agent
agent = create_agent(model=model, tools=tools)

# Usage
result = agent.invoke({"messages": [("human", "Help me draw a shark flying in the sky")]})
      \end{lstlisting}
\end{frame}
%%%%%%%%%%%%%%%%%%%%%%%%%%%%%%%%%%%%%%%%%%%%%%%%%%%%%%%%%%%%%%%%%%%%%%%%%%%%%%%%
\begin{frame}[fragile]\frametitle{Analysis of Results}

      \begin{itemize}
        \item the agent will first understand the task requirements and then decide to use the image tool.
		\item It will write a detailed image description and then call the tool to create the image.
        \item Finally, it will return the generated image URL along with a brief description.
		\item The output might look like this:
      \end{itemize}



      \begin{lstlisting}[language=Python]
{
    "messages": [
        HumanMessage(content='Help me draw a picture of a shark flying in the sky'),
        AIMessage(content='', additional_kwargs={'tool_calls': [...]}),
        ToolMessage(content='https://images.example.com/shark.png'),
        AIMessage(content='Here is your image of a shark flying in the sky: https://images.example.com/shark.png')
    ]
}
      \end{lstlisting}
\end{frame}
